% =====================================================================
% sec/4_experiments.tex  --  Experiments
% =====================================================================

\section{Experiments}
\label{sec:experiments}

\subsection{Datasets}
\label{subsec:datasets}

We use five publicly available colonoscopy datasets:
\begin{itemize}
    \item \textbf{Kvasir-SEG}~\cite{jha2020kvasir}: 1000 polyp images with
          pixel-level annotations. Used for training and in-domain evaluation.
    \item \textbf{CVC-ClinicDB}~\cite{bernal2015clinicdb}: 612 frames from
          31 colonoscopy sequences. Zero-shot evaluation target.
    \item \textbf{ColonDB}~\cite{tajbakhsh2015colondb}: 380 frames.
          Zero-shot evaluation target.
    \item \textbf{CVC-300}~\cite{vazquez2017cvc300}: 60 frames.
          Zero-shot evaluation target.
    \item \textbf{BKAI-IGH}~\cite{ngoc2021bkai}: 1000 frames including
          neoplastic and non-neoplastic polyps. Zero-shot evaluation target.
\end{itemize}

\subsection{Implementation Details}
\label{subsec:impl}

All images are resized to $224{\times}224$. We train for 100 epochs with batch
size 16 using the Adam optimizer ($\text{lr}{=}10^{-4}$, cosine annealing to
$10^{-4}$). The EndoViT backbone is frozen; only the LoRA parameters, text
projection, boundary gate scalars, and seg head are trained
($\approx 3.2$M trainable parameters). Experiments run on a single NVIDIA A100-80GB.

\subsection{Evaluation Metrics}
\label{subsec:metrics}

We report Dice Similarity Coefficient (DSC, \%), Normalized Surface Dice
(NSD, \%) at 2px tolerance, and 95th-percentile Hausdorff Distance (HD95, px).

\subsection{Architecture Ablation (Table~A)}
\label{subsec:ablation}

\cref{tab:ablation} reports the incremental contribution of each proposed
component on the Kvasir-SEG test set (ablation models trained for 30 epochs).

\begin{table}[h]
\centering
\caption{Architecture ablation on Kvasir-SEG (in-domain).}
\label{tab:ablation}
\resizebox{\linewidth}{!}{%
\begin{tabular}{lccccc}
\toprule
\textbf{Model} & \textbf{EndoViT} & \textbf{LoRA} & \textbf{BdFB} & \textbf{DSC\%} & \textbf{NSD\%} \\
\midrule
Baseline (MedCLIPSeg) & \texttimes & \texttimes & \texttimes & -- & -- \\
+ EndoViT             & \checkmark & \texttimes & \texttimes & -- & -- \\
+ LoRA PVL            & \checkmark & \checkmark & \texttimes & -- & -- \\
+ Boundary Feedback   & \checkmark & \checkmark & \checkmark & -- & -- \\
+ TTA (ours)          & \checkmark & \checkmark & \checkmark & \textbf{--} & \textbf{--} \\
\bottomrule
\end{tabular}}
\end{table}

\noindent\textit{(Numeric results will be filled after the Kaggle run completes.)}

\subsection{Cross-Domain Generalization (Table~B)}
\label{subsec:cross_domain}

\cref{tab:cross_domain} reports zero-shot performance of the fully-trained ProbVLP
model on all five datasets, with and without TTA.

\begin{table}[h]
\centering
\caption{Cross-domain generalization results. Model trained on Kvasir-SEG only.}
\label{tab:cross_domain}
\resizebox{\linewidth}{!}{%
\begin{tabular}{lcccccc}
\toprule
\multirow{2}{*}{\textbf{Dataset}} & \multicolumn{3}{c}{\textbf{w/o TTA}} & \multicolumn{3}{c}{\textbf{w/ TTA}} \\
\cmidrule(lr){2-4}\cmidrule(lr){5-7}
 & DSC\% & NSD\% & HD95 & DSC\% & NSD\% & HD95 \\
\midrule
Kvasir (ID)  & -- & -- & -- & -- & -- & -- \\
ClinicDB     & -- & -- & -- & -- & -- & -- \\
ColonDB      & -- & -- & -- & -- & -- & -- \\
CVC-300      & -- & -- & -- & -- & -- & -- \\
BKAI         & -- & -- & -- & -- & -- & -- \\
\midrule
\textit{OOD Avg.} & -- & -- & -- & -- & -- & -- \\
\bottomrule
\end{tabular}}
\end{table}

\noindent\textit{(Numeric results will be filled after the Kaggle run completes.)}
